\begin{problem}[Legacy AG3 Problem 12: Universal property of affine schemes]\label{prob:vi19-p01}
Let $(X,\OO_X)$ be a locally ringed space and let $A$ be a ring.  Prove the natural bijection
\[
\boxed{\operatorname{Hom}_{\mathrm{LRS}}(X,\Spec A)\cong
\operatorname{Hom}_{\mathrm{Ring}}(A,\Gamma(X,\OO_X)).}
\]
Your proof must construct the point map from a ring homomorphism $A\to\Gamma(X,\OO_X)$, prove continuity using principal opens, construct the sheaf morphism by localization, verify locality on stalks, prove uniqueness, and check that the two constructions are inverse.
\end{problem}

\ifdefined\FullProblemDossiers
\paragraph{Why this problem matters.}
This is the universal mapping property that makes affine geometry algebraic.  It contains the affine anti-equivalence as the special case $X=\Spec B$ and is the unique numbered legacy Problem assigned to VI/19.

\paragraph{Useful ideas before starting.}
Given $\alpha:A\to\Gamma(X,\OO_X)$, send $x$ to the inverse image of the maximal ideal of $\OO_{X,x}$ under the germ map.  A section is invertible exactly where all of its germs are units.

\paragraph{Guided hints.}
\begin{enumerate}
\item For $x\in X$, define $\mfp_x=\{a: \alpha(a)_x\in\mfm_x\}$ and prove it is prime.
\item Show $f^{-1}(D(a))$ is the unit locus of the global section $\alpha(a)$.
\item On $D(a)$ localize $A$ and use the inverse of $\alpha(a)$ on the preimage.
\item Check that the stalk map $A_{\mfp_x}\to\OO_{X,x}$ is local.
\item For uniqueness, recover $f(x)$ from the global pullback and compare the sheaf maps on principal opens.
\end{enumerate}

\begin{solution}
First let $f:X\to\Spec A$ be a morphism of locally ringed spaces.  Taking global sections of
$f^\sharp$ gives
\[
\Phi(f):A=\Gamma(\Spec A,\OO)\longrightarrow\Gamma(X,\OO_X).
\]

Conversely, let $\alpha:A\to\Gamma(X,\OO_X)$.  For $x\in X$, compose $\alpha$ with the germ map
$\Gamma(X,\OO_X)\to\OO_{X,x}$.  Since $\OO_{X,x}$ is local with maximal ideal $\mfm_x$, put
\[
f_\alpha(x)=\{a\in A:\alpha(a)_x\in\mfm_x\}.
\]
The inverse image of a prime ideal under a ring homomorphism is prime, so $f_\alpha(x)\in\Spec A$.

For $a\in A$,
\[
x\in f_\alpha^{-1}(D(a))
\iff \alpha(a)_x\notin\mfm_x
\iff \alpha(a)_x\in\OO_{X,x}^{\times}.
\]
The unit locus of a section is open: if a germ is invertible, representatives of the section and its germwise inverse multiply to $1$ on some neighborhood.  Hence $f_\alpha^{-1}(D(a))$ is open.  Principal opens form a basis, so $f_\alpha$ is continuous.

On $D(a)\subseteq\Spec A$, the section ring is $A_a$.  Over $U_a=f_\alpha^{-1}(D(a))$, the section $\alpha(a)|_{U_a}$ is a unit.  Therefore the universal property of localization gives
\[
A_a\longrightarrow\Gamma(U_a,\OO_X),\qquad
\frac{r}{a^n}\longmapsto
\alpha(r)|_{U_a}\,\alpha(a)|_{U_a}^{-n}.
\]
For $D(ab)\subseteq D(a)$ these maps commute with restriction, so they define a morphism on the principal-open basis and hence a sheaf morphism
\[
f_\alpha^\sharp:\OO_{\Spec A}\longrightarrow (f_\alpha)_*\OO_X.
\]
At $x$, with $\mfp=f_\alpha(x)$, the induced stalk map is
\[
A_\mfp\longrightarrow\OO_{X,x},\qquad a/s\longmapsto\alpha(a)_x\alpha(s)_x^{-1}.
\]
By definition, $a$ maps into $\mfm_x$ exactly when $a\in\mfp$, so the inverse image of $\mfm_x$ is $\mfp A_\mfp$.  The stalk map is local.  Hence $(f_\alpha,f_\alpha^\sharp)$ is a morphism of locally ringed spaces.

To prove uniqueness, suppose $g:X\to\Spec A$ has global pullback $\alpha$.  Locality of the stalk map gives
\[
g(x)=\{a\in A:\alpha(a)_x\in\mfm_x\}=f_\alpha(x),
\]
so $g$ and $f_\alpha$ have the same point map.  On every principal open $D(a)$, the sheaf map must send $r/a^n$ to $\alpha(r)\alpha(a)^{-n}$, so it also equals $f_\alpha^\sharp$.  Thus $g=f_\alpha$.

Starting with $\alpha$, the constructed morphism has global pullback $\alpha$ by taking $a=1$.  Starting with $f$, the reconstruction from its global pullback recovers both its point map and its sheaf map by the preceding uniqueness argument.  Hence the constructions are inverse.  They are natural because restriction, localization, and composition of pullbacks are functorial.
\end{solution}

\paragraph{What happened mathematically.}
The target affine scheme is completely controlled by its global ring: a map into it is exactly a compatible way to pull every global affine coordinate back to the source.

\paragraph{Extensions.}
Specialize to $X=\Spec B$ to derive the affine anti-equivalence; specialize to $A=\mathbb Z[t]$ to identify global sections with maps to the affine line.

\paragraph{Provenance.}
\textbf{Legacy recovery: theory-of-algebraic-geometry-3.tex, Problem 12. Preserved at full universal-property depth and expanded without changing its mathematical center.}
\else
\paragraph{Provenance.} \textbf{Legacy recovery: theory-of-algebraic-geometry-3.tex, Problem 12. Preserved at full universal-property depth and expanded without changing its mathematical center.}
\fi
